Comorbidity networks provide a powerful framework for understanding disease relationships beyond simple co-occurrence statistics, revealing shared physiological mechanisms and guiding clinical management of multimorbidity. Understanding multimorbidity patterns is essential for developing effective, integrated, and person-centered care models~\cite{Chiadi2023}. A crucial advancement lies in analyzing disease combinations that exceed what would be expected by chance, addressing the high variability in how diseases co-occur, moving beyond simple prevalence rates and introducing patient-specific characteristics as risk factors ~\cite{EPF2012ChronicDiseaseResponse}. This shift is particularly important because knowledge of multimorbidity patterns has direct implications for patient-oriented (rather than disease-oriented) prevention, diagnosis, treatment, and prognosis, and, as highlighted by the National Institute for Health and Clinical Excellence, provides essential information for developing clinical guidelines and decision-support tools for managing patients with multiple chronic conditions~\cite{NICE2016NG56}. The integration of temporal information from longitudinal patient sequences represents a critical methodological advancement in this direction. By capturing both the dynamics and temporal ordering of disease progression, sequence-based methods can distinguish between early- and late-emerging conditions in the multimorbidity trajectory and identify the order in which diseases develop. Incorporating this temporal information allows for the construction of a more “realistic” comorbidity network, improving our understanding of the relationships between co-occurring diseases and informing the design of targeted intervention strategies~\cite{siggaard2020disease}. 

We developed a comprehensive methodological pipeline for constructing disease comorbidity from longitudinal patient data designed to overcome three key limitations of existing comorbidity network approaches: the absence of temporal information, confounding, and reliance on arbitrary thresholding for sparsification. Our approach adapts Word2Vec embeddings to learn disease representations from longitudinal patient trajectories, treating diseases as "words" and patient sequences as "sentences". This embedding-based framework offers several advantages over traditional co-occurrence statistics: dense 200-dimensional vectors capture richer information than scalar similarity measures, encoding both semantic relationships and temporal context; the continuous embedding space enables gradient-based optimization and transfer learning from large biomedical corpora; and incorporating risk factors allows shared patterns to be represented directly in the vector space, improving similarity estimates by reducing spurious associations driven by common exposures. Building on this principle, we introduce a confounding adjustment framework by jointly embedding diseases and risk factors. The model redistributes association signals: correlations driven by shared risk factors are allocated to risk factor dimensions, yielding adjusted disease embeddings that better reflect direct pathophysiological relationships. The difference between risk-adjusted and unadjusted similarities provides an interpretable measure of confounding magnitude for each disease pair, validated through a three-step framework demonstrating that observed decreases reflect genuine confounding removal rather than random variation. While we cannot definitively establish causality, the systematic correlation between similarity changes and shared risk factors provides strong evidence that our adjustment procedure successfully removes confounding. Applied to UK Biobank data with 24 cardiometabolic diseases and 76 risk factors, our framework revealed substantial confounding effects: 91.3\% of disease pairs showed decreased similarity after risk adjustment, with highly affected pairs systematically sharing more risk factors. Hypertension and diabetes mellitus appeared in 80\% of most-affected pairs, confirming that their frequent co-occurrence in unadjusted networks is heavily driven by shared lifestyle factors and comorbidities. Conversely, acute clinical events (acute myocardial infarction, cardiac arrest) and advanced interventions (heart transplantation, pacemaker implantation) showed minimal changes, appropriately reflecting their nature as acute events or end-stage treatments less confounded by gradual risk factor accumulation. For network sparsification, Gaussian Graphical Models provided a statistically principled alternative to traditional backboning methods~\cite{serrano2009extracting}, which fail when similarity values are homogeneous. GGM naturally handles such cases by modeling the full covariance structure and identifying edges that remain significant after conditioning on the entire network. Topological analysis of the GGM-derived network identified atrial fibrillation, acute kidney injury, and coronary disease as hub diseases with critical bridging functions, while community detection recovered three clinically meaningful modules (cardiorenal-metabolic, primary cardiovascular, large vessel cerebrovascular) aligning with established cardiometabolic taxonomy~(\cite{Chiadi2023, WHO_CVDs_2025}).

Beyond methodological contributions, our framework enables several clinically relevant applications that advance evidence-based practice in multimorbidity management. First, the topology-based identification of hub diseases and bridging nodes provides actionable insights for prioritizing preventive interventions. Diseases with high betweenness centrality, such as atrial fibrillation and acute kidney injury, represent strategic targets where interventions may have broad systemic impact across multiple pathophysiological domains~\cite{Zhao2023ComorbidityNetwork}. Second, recent research has shown that comorbidity networks can evolve over time, with disease association patterns shifting across life stages and phases of progression~(\citealp{dervic2025comorbidity, gardinazzi2025characterization}). Our embedding-based framework can be readily extended to construct age-specific networks by applying the pipeline to stratified cohorts based on age. This reveal how multimorbidity patterns change over the life course and how the influence of risk factors varies with age. Methodologically, temporal differences can be identified in the same way we compared adjusted and unadjusted networks, by systematically contrasting networks constructed under different modeling specifications. Third, patient-specific representations create new opportunities for personalized risk stratification. By mapping individual patient trajectories into the learned embedding space, representing each patient as the mean embedding of the diseases recorded in their clinical history, we can cluster patients with similar comorbidity progression patterns. This enables stratified survival analyses and the identification of profiles that deviate from typical trajectories, thereby supporting earlier detection and intervention for high-risk multimorbidity phenotypes~\cite{liang2024comorbidity}.Fourth, multimorbidity patterns identified through our framework can serve as predictors of intense healthcare usage and adverse health outcomes. As demonstrated by~\cite{Wolff2002}, the simultaneous presentation of diseases may have multiplicative rather than additive effects on health and healthcare costs. Network-derived features, such as membership in high-risk disease modules or proximity to hub diseases, could enhance prediction models for disability, adverse drug events, and increased disease severity, ultimately supporting resource allocation and care planning decisions~\cite{buja2025healthcare}.

Our study has several limitations that suggest directions for methodological development. First, while we adjust for 76 measured risk factors, residual confounding from unmeasured variables may persist. Future work could enrich the set of information by integrating genetic, proteomic, and metabolomic data to reveal shared underlying biological mechanisms and identify potential therapeutic targets that extend across multiple conditions~\citep{Sibilio2025, nam2020network}. Second, the embedding approach assumes that diseases appearing in similar temporal contexts have similar meanings, which may not hold for all disease pairs. Rare diseases with few occurrences may be poorly represented in the embedding space, potentially missing important but uncommon comorbidity patterns. Additionally, diseases with multiple distinct etiologies (e.g., heart failure with preserved or reduced ejection fraction) are treated as single entities, potentially obscuring heterogeneity. Future extensions could incorporate disease subtypes when data availability and sufficient sample sizes permit. Third, while we capture temporal ordering at the disease level, we do not explicitly model time intervals between events or disease onset ages, which may contain prognostic information. In addition, the estimated similarity networks are undirected, precluding inference of disease progression directionality or causal pathways. Although temporal sequence information informs the embeddings, the resulting networks represent associations rather than directional dependencies. 
Incorporating causal inference frameworks~(\citealp{pearl2009causality, prosperi2020causal, xu2019patient, kong_causal}) could enable directed network estimation, distinguishing causes from effects in disease relationships and supporting the construction of age-dependent or temporally evolving networks that track disease progression trajectories over the life course.

In conclusion, this work demonstrates that integrating longitudinal patient data with explicit risk factor adjustment and advanced representation learning substantially enhances the validity and clinical utility of disease network analyses. By addressing key limitations in existing approaches our framework provides a statistically rigorous foundation for characterizing multimorbidity patterns. The systematic pipeline we present offers a generalizable strategy applicable beyond cardiometabolic conditions, with direct implications for evidence-based guideline development, patient stratification and healthcare resource planning. These applications link research with practice, paving the way for more precise characterization of multimorbidity patterns and their underlying mechanisms across diverse disease domains.


%%COMMENTI CLINICI RIMOSSI
%Highly affected disease pairs shared substantially more risk factors than lower-ranked pairs, with all 14 most-affected pairs showing decreased similarity after adjustment. Similarly, associations involving mitral and tricuspid valve insufficiency are highly confounded by shared hemodynamic risks such as hypotension, cardiac complications, and conduction disorders, inflating their apparent relationships with metabolic conditions in unadjusted networks. Collectively, our findings suggest that the network topology effectively recovers clinical progression pathways, identifying strategic nodes governing transitions between distinct pathological states. In particular, the dual prominence of atrial fibrillation highlights its high local connectivity and strategic bridging function, identifying it as a key mechanistic hub where interventions may have broad systemic impact. Community detection identified three disease modules, metabolic, primary cardiovascular, and large vessel cerebrovascular, aligning with established cardiometabolic disease taxonomy and demonstrating that our data-driven approach naturally recovers clinically meaningful disease groupings~(\citealp{Chiadi2023, WHO_CVDs_2025}).